% =====================================================================
%  MSACL DS301 Deep Learning · Quiz 10 · Imbalanced Data
%  Academic problem-set style (single-source, \ifsolution toggle)
%  SINGLE SOURCE — produces BOTH the student quiz and the answer key.
%
%  ANSWER TOGGLE
%  -------------
%  Student version (answers HIDDEN) — the default:
%       pdflatex quiz10_imbalanced_data.tex
%  Answer key (answers SHOWN):
%       pdflatex "\def\showsolutions{}% =====================================================================
%  MSACL DS301 Deep Learning · Quiz 10 · Imbalanced Data
%  Academic problem-set style (single-source, \ifsolution toggle)
%  SINGLE SOURCE — produces BOTH the student quiz and the answer key.
%
%  ANSWER TOGGLE
%  -------------
%  Student version (answers HIDDEN) — the default:
%       pdflatex quiz10_imbalanced_data.tex
%  Answer key (answers SHOWN):
%       pdflatex "\def\showsolutions{}% =====================================================================
%  MSACL DS301 Deep Learning · Quiz 10 · Imbalanced Data
%  Academic problem-set style (single-source, \ifsolution toggle)
%  SINGLE SOURCE — produces BOTH the student quiz and the answer key.
%
%  ANSWER TOGGLE
%  -------------
%  Student version (answers HIDDEN) — the default:
%       pdflatex quiz10_imbalanced_data.tex
%  Answer key (answers SHOWN):
%       pdflatex "\def\showsolutions{}% =====================================================================
%  MSACL DS301 Deep Learning · Quiz 10 · Imbalanced Data
%  Academic problem-set style (single-source, \ifsolution toggle)
%  SINGLE SOURCE — produces BOTH the student quiz and the answer key.
%
%  ANSWER TOGGLE
%  -------------
%  Student version (answers HIDDEN) — the default:
%       pdflatex quiz10_imbalanced_data.tex
%  Answer key (answers SHOWN):
%       pdflatex "\def\showsolutions{}\input{quiz10_imbalanced_data.tex}"
%  (or, to flip by hand, change \solutionfalse to \solutiontrue below.)
% =====================================================================
\documentclass[11pt]{article}

\newif\ifsolution
\ifdefined\showsolutions \solutiontrue \else \solutionfalse \fi
% \solutiontrue   % <-- uncomment to hard-code the ANSWER KEY build

\usepackage[letterpaper, margin=0.6in, top=0.7in, bottom=0.5in]{geometry}
\usepackage{amsmath, amssymb}
\usepackage[table]{xcolor}
\usepackage{array}
\usepackage{fancyhdr}

\pagestyle{fancy}
\fancyhf{}
\fancyhead[L]{\small\textbf{MSACL DS301 --- Deep Learning}}
\fancyhead[R]{\small\textbf{Quiz 10: Imbalanced Data\ifsolution\ \textmd{(Answer Key)}\fi}}
\fancyfoot[C]{\thepage}
\renewcommand{\headrulewidth}{0.4pt}

\newcommand{\blk}[1]{\underline{\hspace{#1}}}
% Reveal-or-blank: shows the answer in the key, a blank rule on the student sheet.
\newcommand{\rb}[2]{\ifsolution\textbf{#1}\else\blk{#2}\fi}
% Boxed answer (key only) and blank writing space (student only).
\newcommand{\keybox}[1]{\ifsolution\par\vspace{2pt}%
  \fbox{\parbox{0.965\textwidth}{\small\textbf{Answer.}\; #1}}\par\fi}
\newcommand{\work}[1]{\ifsolution\else\vspace{#1}\fi}

% ---- course palette (numeric grids) --------------------------------------
\definecolor{ink}{HTML}{1B2025}
\definecolor{teal}{HTML}{0E7C7B}
\definecolor{tealsoft}{HTML}{E3F0EF}
\definecolor{amber}{HTML}{E09E2F}
\definecolor{ambersoft}{HTML}{FBF0DC}
\definecolor{redd}{HTML}{C4534F}
\definecolor{redsoft}{HTML}{F7E4E3}
\definecolor{inksoft}{HTML}{3D444C}
\definecolor{muted}{HTML}{8A9099}
\definecolor{hair}{HTML}{D9D9D2}
% square, centred cells for the confusion matrix
\newcolumntype{G}{>{\centering\arraybackslash}p{28mm}}
\newcommand{\gr}{\renewcommand{\arraystretch}{1.7}}

\setlength{\parindent}{0pt}
\setlength{\parskip}{3pt}
\setlength{\abovedisplayskip}{5pt}
\setlength{\belowdisplayskip}{5pt}

\begin{document}

\begin{center}
{\Large \textbf{Quiz 10 --- Imbalanced Data}}\\[2pt]
{\normalsize Compute precision / recall / specificity, choose AUROC vs.\ AUPRC, match a treatment, and interpolate a SMOTE point}
\end{center}

\vspace{-2pt}
\begin{center}
\fbox{\parbox{0.92\textwidth}{\small
\textbf{Instructions:}\; Answer all four questions --- every part is answerable from the formulas and pictures on
the slides. A rapid resistance screen was run on \textbf{500 isolates} (positive $=$ \emph{resistant}); Q1--Q2 use
its confusion matrix. Compute each rate exactly; no calculators needed.%
\ifsolution\else\\[4pt]\textbf{Name:}\ \blk{6cm}\hfill\textbf{Date:}\ \blk{3.5cm}\fi
}}
\end{center}

\vspace{-1pt}
\begin{center}\small
\textbf{Formulas (from the slides):}\;
$\text{recall (sensitivity)} = \dfrac{TP}{TP+FN}$ \quad
$\text{specificity} = \dfrac{TN}{TN+FP}$ \quad
$\text{precision} = \dfrac{TP}{TP+FP}$ \quad
$\text{FPR} = \dfrac{FP}{FP+TN}$
\end{center}

\vspace{-2pt}\hrule\vspace{5pt}

% =========================== Q1 ===========================
\textbf{1. Three rates by hand.}\; {\small The screen's confusion matrix (rows $=$ the model's call, columns $=$
truth):}

\vspace{4pt}
\begin{center}
\gr\begin{tabular}{r|G|G|}
\multicolumn{1}{r}{} & \multicolumn{1}{c}{\small truth: $+$ (resistant)} & \multicolumn{1}{c}{\small truth: $-$ (susceptible)}\\\cline{2-3}
\small call: $+$ & \cellcolor{tealsoft}TP $=$ 16 & \cellcolor{ambersoft}FP $=$ 32 \\\cline{2-3}
\small call: $-$ & \cellcolor{redsoft}FN $=$ 4 & \cellcolor{tealsoft}TN $=$ 448 \\\cline{2-3}
\end{tabular}\\[3pt]
{\footnotesize 20 resistant \;$\cdot$\; 480 susceptible \;$\cdot$\; 500 total}
\end{center}

\vspace{2pt}
\textbf{(a)} recall (sensitivity) $=$ \rb{$16/20 = 0.80$ \,(80\%)}{3.4cm}
\work{0.3cm}

\textbf{(b)} specificity $=$ \rb{$448/480 = 0.93$ \,(93\%)}{3.4cm}
\work{0.3cm}

\textbf{(c)} precision $=$ \rb{$16/48 = 0.33$ \,(33\%)}{3.4cm}
\work{0.3cm}

\textbf{(d)} A model that ALWAYS calls ``susceptible'' scores accuracy $= 480/500 = 96\%$. The lab must not miss
resistant isolates. \emph{Which rate(s) do you trust here, and why is accuracy misleading?}

\work{1.0cm}
\ifsolution\else\blk{\linewidth}\par\vspace{2pt}\blk{\linewidth}\fi

\keybox{(a) recall $= 16/(16+4) = 16/20 = \mathbf{0.80}$ (\textbf{80\%}). (b) specificity $= 448/(448+32) = 448/480 =
\mathbf{0.93}$ (\textbf{93\%}). (c) precision $= 16/(16+32) = 16/48 = \mathbf{0.33}$ (\textbf{33\%}). Accuracy $=
(16+448)/500 = 464/500 = \mathbf{0.93}$ (93\%). (d) Trust \textbf{recall} (catch the resistant) reported \emph{with}
\textbf{precision} (only 1 flag in 3 is real). Positives are rare (20 of 500), so accuracy is dominated by the 480
susceptibles: an always-``susceptible'' model scores \textbf{96\%} yet catches \textbf{0 of 20} resistant. Accuracy
hides both the missed cases and the false alarms.}

\vspace{2pt}\hrule\vspace{5pt}

% =========================== Q2 ===========================
\textbf{2. AUROC or AUPRC?}\; {\small Same screen. Positives are \textbf{rare: 20 of 500 $= 4\%$ prevalence.} Plot
this one operating point on each curve, then choose what to report across all thresholds.}

\vspace{2pt}
\textbf{(a)} ROC point: FPR $=\ \dfrac{32}{32+448} =\ $ \rb{$32/480 = 0.067$}{2.8cm}, \; TPR (recall) $=\ $
\rb{$0.80$}{2.0cm}\; --- does it hug the top-left corner? \rb{yes (looks great)}{3.0cm}
\work{0.3cm}

\textbf{(b)} PR point: recall $=\ $ \rb{$0.80$}{2.0cm}, \; precision $=\ $ \rb{$0.33$}{2.0cm};\;
the PR chance baseline sits at the prevalence $=\ $ \rb{$20/500 = 0.04$}{2.8cm}
\work{0.3cm}

\textbf{(c)} To report this screen across \emph{all thresholds}, would you trust the \textbf{AUROC} or the
\textbf{AUPRC / PR curve}, and why?

\work{1.0cm}
\ifsolution\else\blk{\linewidth}\par\vspace{2pt}\blk{\linewidth}\fi

\keybox{(a) FPR $= 32/(32+448) = 32/480 = \mathbf{0.067}$; TPR $= \mathbf{0.80}$ --- the point $(0.067,\,0.80)$
\textbf{hugs the top-left}, so the ROC/AUROC \emph{looks excellent}. (b) PR point $=(\text{recall }\mathbf{0.80},\,
\text{precision }\mathbf{0.33})$; PR baseline $=$ prevalence $= 20/500 = \mathbf{0.04}$ (ROC's baseline is $0.5$).
(c) Trust the \textbf{AUPRC / PR curve} (reported with recall). At 4\% prevalence the 480 true-negatives dominate, so a
tiny FPR keeps AUROC flatteringly high even though precision is only \textbf{0.33} (2 of 3 flags false). AUROC is
invariant to the class ratio; AUPRC collapses toward the prevalence, so it stays honest as recall grows
(Saito \& Rehmsmeier 2015).}

\vspace{2pt}\hrule\vspace{5pt}

% =========================== Q3 ===========================
\textbf{3. Match the treatment.}\; {\small Match each situation to \emph{one} move from the menu:
\textbf{curation} \;$\cdot$\; \textbf{class weights / SMOTE / oversample} \;$\cdot$\; \textbf{focal loss}
\;$\cdot$\; \textbf{curriculum}. Each is used once.}

\vspace{2pt}
\textbf{(a)} Junk runs and mislabeled spectra contaminate the training set --- and the same isolate
appears three times. $\rightarrow$\ \rb{curation}{5cm}
\work{0.25cm}
\textbf{(b)} Only 41 real resistant spectra; the model just predicts ``susceptible''. $\rightarrow$\
\rb{class weights / SMOTE / oversample}{5cm}
\work{0.25cm}
\textbf{(c)} Easy cases are fine, but borderline isolates are missed. $\rightarrow$\ \rb{focal loss}{5cm}
\work{0.25cm}
\textbf{(d)} Training is unstable when very easy and very hard batches alternate. $\rightarrow$\
\rb{curriculum}{5cm}
\work{0.35cm}

\textbf{(e)} \emph{Which of these four adds NEW resistant signal to the dataset?}

\work{0.9cm}
\ifsolution\else\blk{\linewidth}\par\fi

\keybox{(a) \textbf{curation} --- clean the labels and drop the junk before anything fancier, then
\emph{deduplicate}: embed every spectrum (the CNN's last hidden layer) and treat $\cos \geq 0.99$ as the
same isolate, keeping one --- repeats split across train and test are leakage. (b) \textbf{class weights / SMOTE /
oversample} --- make the rare class count. (c) \textbf{focal loss} --- up-weight the hard, low-confidence cases.
(d) \textbf{curriculum} --- order easy$\to$hard, or class-balanced$\to$true prevalence. (e) \textbf{None of them.}
Every move re-uses the same 41 resistant spectra (reweighting shouts them louder, SMOTE interpolates between them,
focal loss focuses on them, curriculum reorders them) --- \emph{no new biological information enters}. Only
\textbf{collecting more real resistant isolates} adds signal.}

\vspace{2pt}\hrule\vspace{5pt}

% =========================== Q4 ===========================
\textbf{4. SMOTE by hand.}\; {\small SMOTE makes a new minority point on the segment between two real ones:
\; $\text{synthetic} = A + \lambda\,(B-A)$. Use two $m/z$-bin intensities $A=(2,\,6)$ and $B=(4,\,10)$.}

\vspace{2pt}
\textbf{(a)} With $\lambda = 0.25$:\; synthetic $=\ (\,2 + 0.25\,(4-2),\ \ 6 + 0.25\,(10-6)\,) =\ $
\rb{$(2.5,\ 7)$}{3cm}
\work{0.4cm}

\textbf{(b)} \emph{One sentence: why is raw-spectrum SMOTE risky in 6{,}000 dimensions?}

\work{1.0cm}
\ifsolution\else\blk{\linewidth}\par\vspace{2pt}\blk{\linewidth}\fi

\keybox{(a) synthetic $= (2 + 0.25\cdot 2,\ 6 + 0.25\cdot 4) = (2 + 0.5,\ 6 + 1) = \mathbf{(2.5,\ 7)}$ --- a quarter
of the way from $A$ toward $B$. (b) A straight-line blend of two real 6{,}000-dimensional spectra can produce a
\textbf{non-physical} spectrum (a peak pattern no instrument would ever record), so pair SMOTE with \emph{curation}
and domain augmentation ($m/z$ jitter), or simply \emph{oversample real} spectra.}

\vspace{4pt}\hrule\vspace{3pt}
{\small\itshape \ifsolution Answer key --- other defensible wording earns credit if the numbers match.%
\else Keep this sheet --- it's your checklist for any rare-class assay in your own lab.\fi}

\end{document}
"
%  (or, to flip by hand, change \solutionfalse to \solutiontrue below.)
% =====================================================================
\documentclass[11pt]{article}

\newif\ifsolution
\ifdefined\showsolutions \solutiontrue \else \solutionfalse \fi
% \solutiontrue   % <-- uncomment to hard-code the ANSWER KEY build

\usepackage[letterpaper, margin=0.6in, top=0.7in, bottom=0.5in]{geometry}
\usepackage{amsmath, amssymb}
\usepackage[table]{xcolor}
\usepackage{array}
\usepackage{fancyhdr}

\pagestyle{fancy}
\fancyhf{}
\fancyhead[L]{\small\textbf{MSACL DS301 --- Deep Learning}}
\fancyhead[R]{\small\textbf{Quiz 10: Imbalanced Data\ifsolution\ \textmd{(Answer Key)}\fi}}
\fancyfoot[C]{\thepage}
\renewcommand{\headrulewidth}{0.4pt}

\newcommand{\blk}[1]{\underline{\hspace{#1}}}
% Reveal-or-blank: shows the answer in the key, a blank rule on the student sheet.
\newcommand{\rb}[2]{\ifsolution\textbf{#1}\else\blk{#2}\fi}
% Boxed answer (key only) and blank writing space (student only).
\newcommand{\keybox}[1]{\ifsolution\par\vspace{2pt}%
  \fbox{\parbox{0.965\textwidth}{\small\textbf{Answer.}\; #1}}\par\fi}
\newcommand{\work}[1]{\ifsolution\else\vspace{#1}\fi}

% ---- course palette (numeric grids) --------------------------------------
\definecolor{ink}{HTML}{1B2025}
\definecolor{teal}{HTML}{0E7C7B}
\definecolor{tealsoft}{HTML}{E3F0EF}
\definecolor{amber}{HTML}{E09E2F}
\definecolor{ambersoft}{HTML}{FBF0DC}
\definecolor{redd}{HTML}{C4534F}
\definecolor{redsoft}{HTML}{F7E4E3}
\definecolor{inksoft}{HTML}{3D444C}
\definecolor{muted}{HTML}{8A9099}
\definecolor{hair}{HTML}{D9D9D2}
% square, centred cells for the confusion matrix
\newcolumntype{G}{>{\centering\arraybackslash}p{28mm}}
\newcommand{\gr}{\renewcommand{\arraystretch}{1.7}}

\setlength{\parindent}{0pt}
\setlength{\parskip}{3pt}
\setlength{\abovedisplayskip}{5pt}
\setlength{\belowdisplayskip}{5pt}

\begin{document}

\begin{center}
{\Large \textbf{Quiz 10 --- Imbalanced Data}}\\[2pt]
{\normalsize Compute precision / recall / specificity, choose AUROC vs.\ AUPRC, match a treatment, and interpolate a SMOTE point}
\end{center}

\vspace{-2pt}
\begin{center}
\fbox{\parbox{0.92\textwidth}{\small
\textbf{Instructions:}\; Answer all four questions --- every part is answerable from the formulas and pictures on
the slides. A rapid resistance screen was run on \textbf{500 isolates} (positive $=$ \emph{resistant}); Q1--Q2 use
its confusion matrix. Compute each rate exactly; no calculators needed.%
\ifsolution\else\\[4pt]\textbf{Name:}\ \blk{6cm}\hfill\textbf{Date:}\ \blk{3.5cm}\fi
}}
\end{center}

\vspace{-1pt}
\begin{center}\small
\textbf{Formulas (from the slides):}\;
$\text{recall (sensitivity)} = \dfrac{TP}{TP+FN}$ \quad
$\text{specificity} = \dfrac{TN}{TN+FP}$ \quad
$\text{precision} = \dfrac{TP}{TP+FP}$ \quad
$\text{FPR} = \dfrac{FP}{FP+TN}$
\end{center}

\vspace{-2pt}\hrule\vspace{5pt}

% =========================== Q1 ===========================
\textbf{1. Three rates by hand.}\; {\small The screen's confusion matrix (rows $=$ the model's call, columns $=$
truth):}

\vspace{4pt}
\begin{center}
\gr\begin{tabular}{r|G|G|}
\multicolumn{1}{r}{} & \multicolumn{1}{c}{\small truth: $+$ (resistant)} & \multicolumn{1}{c}{\small truth: $-$ (susceptible)}\\\cline{2-3}
\small call: $+$ & \cellcolor{tealsoft}TP $=$ 16 & \cellcolor{ambersoft}FP $=$ 32 \\\cline{2-3}
\small call: $-$ & \cellcolor{redsoft}FN $=$ 4 & \cellcolor{tealsoft}TN $=$ 448 \\\cline{2-3}
\end{tabular}\\[3pt]
{\footnotesize 20 resistant \;$\cdot$\; 480 susceptible \;$\cdot$\; 500 total}
\end{center}

\vspace{2pt}
\textbf{(a)} recall (sensitivity) $=$ \rb{$16/20 = 0.80$ \,(80\%)}{3.4cm}
\work{0.3cm}

\textbf{(b)} specificity $=$ \rb{$448/480 = 0.93$ \,(93\%)}{3.4cm}
\work{0.3cm}

\textbf{(c)} precision $=$ \rb{$16/48 = 0.33$ \,(33\%)}{3.4cm}
\work{0.3cm}

\textbf{(d)} A model that ALWAYS calls ``susceptible'' scores accuracy $= 480/500 = 96\%$. The lab must not miss
resistant isolates. \emph{Which rate(s) do you trust here, and why is accuracy misleading?}

\work{1.0cm}
\ifsolution\else\blk{\linewidth}\par\vspace{2pt}\blk{\linewidth}\fi

\keybox{(a) recall $= 16/(16+4) = 16/20 = \mathbf{0.80}$ (\textbf{80\%}). (b) specificity $= 448/(448+32) = 448/480 =
\mathbf{0.93}$ (\textbf{93\%}). (c) precision $= 16/(16+32) = 16/48 = \mathbf{0.33}$ (\textbf{33\%}). Accuracy $=
(16+448)/500 = 464/500 = \mathbf{0.93}$ (93\%). (d) Trust \textbf{recall} (catch the resistant) reported \emph{with}
\textbf{precision} (only 1 flag in 3 is real). Positives are rare (20 of 500), so accuracy is dominated by the 480
susceptibles: an always-``susceptible'' model scores \textbf{96\%} yet catches \textbf{0 of 20} resistant. Accuracy
hides both the missed cases and the false alarms.}

\vspace{2pt}\hrule\vspace{5pt}

% =========================== Q2 ===========================
\textbf{2. AUROC or AUPRC?}\; {\small Same screen. Positives are \textbf{rare: 20 of 500 $= 4\%$ prevalence.} Plot
this one operating point on each curve, then choose what to report across all thresholds.}

\vspace{2pt}
\textbf{(a)} ROC point: FPR $=\ \dfrac{32}{32+448} =\ $ \rb{$32/480 = 0.067$}{2.8cm}, \; TPR (recall) $=\ $
\rb{$0.80$}{2.0cm}\; --- does it hug the top-left corner? \rb{yes (looks great)}{3.0cm}
\work{0.3cm}

\textbf{(b)} PR point: recall $=\ $ \rb{$0.80$}{2.0cm}, \; precision $=\ $ \rb{$0.33$}{2.0cm};\;
the PR chance baseline sits at the prevalence $=\ $ \rb{$20/500 = 0.04$}{2.8cm}
\work{0.3cm}

\textbf{(c)} To report this screen across \emph{all thresholds}, would you trust the \textbf{AUROC} or the
\textbf{AUPRC / PR curve}, and why?

\work{1.0cm}
\ifsolution\else\blk{\linewidth}\par\vspace{2pt}\blk{\linewidth}\fi

\keybox{(a) FPR $= 32/(32+448) = 32/480 = \mathbf{0.067}$; TPR $= \mathbf{0.80}$ --- the point $(0.067,\,0.80)$
\textbf{hugs the top-left}, so the ROC/AUROC \emph{looks excellent}. (b) PR point $=(\text{recall }\mathbf{0.80},\,
\text{precision }\mathbf{0.33})$; PR baseline $=$ prevalence $= 20/500 = \mathbf{0.04}$ (ROC's baseline is $0.5$).
(c) Trust the \textbf{AUPRC / PR curve} (reported with recall). At 4\% prevalence the 480 true-negatives dominate, so a
tiny FPR keeps AUROC flatteringly high even though precision is only \textbf{0.33} (2 of 3 flags false). AUROC is
invariant to the class ratio; AUPRC collapses toward the prevalence, so it stays honest as recall grows
(Saito \& Rehmsmeier 2015).}

\vspace{2pt}\hrule\vspace{5pt}

% =========================== Q3 ===========================
\textbf{3. Match the treatment.}\; {\small Match each situation to \emph{one} move from the menu:
\textbf{curation} \;$\cdot$\; \textbf{class weights / SMOTE / oversample} \;$\cdot$\; \textbf{focal loss}
\;$\cdot$\; \textbf{curriculum}. Each is used once.}

\vspace{2pt}
\textbf{(a)} Junk runs and mislabeled spectra contaminate the training set --- and the same isolate
appears three times. $\rightarrow$\ \rb{curation}{5cm}
\work{0.25cm}
\textbf{(b)} Only 41 real resistant spectra; the model just predicts ``susceptible''. $\rightarrow$\
\rb{class weights / SMOTE / oversample}{5cm}
\work{0.25cm}
\textbf{(c)} Easy cases are fine, but borderline isolates are missed. $\rightarrow$\ \rb{focal loss}{5cm}
\work{0.25cm}
\textbf{(d)} Training is unstable when very easy and very hard batches alternate. $\rightarrow$\
\rb{curriculum}{5cm}
\work{0.35cm}

\textbf{(e)} \emph{Which of these four adds NEW resistant signal to the dataset?}

\work{0.9cm}
\ifsolution\else\blk{\linewidth}\par\fi

\keybox{(a) \textbf{curation} --- clean the labels and drop the junk before anything fancier, then
\emph{deduplicate}: embed every spectrum (the CNN's last hidden layer) and treat $\cos \geq 0.99$ as the
same isolate, keeping one --- repeats split across train and test are leakage. (b) \textbf{class weights / SMOTE /
oversample} --- make the rare class count. (c) \textbf{focal loss} --- up-weight the hard, low-confidence cases.
(d) \textbf{curriculum} --- order easy$\to$hard, or class-balanced$\to$true prevalence. (e) \textbf{None of them.}
Every move re-uses the same 41 resistant spectra (reweighting shouts them louder, SMOTE interpolates between them,
focal loss focuses on them, curriculum reorders them) --- \emph{no new biological information enters}. Only
\textbf{collecting more real resistant isolates} adds signal.}

\vspace{2pt}\hrule\vspace{5pt}

% =========================== Q4 ===========================
\textbf{4. SMOTE by hand.}\; {\small SMOTE makes a new minority point on the segment between two real ones:
\; $\text{synthetic} = A + \lambda\,(B-A)$. Use two $m/z$-bin intensities $A=(2,\,6)$ and $B=(4,\,10)$.}

\vspace{2pt}
\textbf{(a)} With $\lambda = 0.25$:\; synthetic $=\ (\,2 + 0.25\,(4-2),\ \ 6 + 0.25\,(10-6)\,) =\ $
\rb{$(2.5,\ 7)$}{3cm}
\work{0.4cm}

\textbf{(b)} \emph{One sentence: why is raw-spectrum SMOTE risky in 6{,}000 dimensions?}

\work{1.0cm}
\ifsolution\else\blk{\linewidth}\par\vspace{2pt}\blk{\linewidth}\fi

\keybox{(a) synthetic $= (2 + 0.25\cdot 2,\ 6 + 0.25\cdot 4) = (2 + 0.5,\ 6 + 1) = \mathbf{(2.5,\ 7)}$ --- a quarter
of the way from $A$ toward $B$. (b) A straight-line blend of two real 6{,}000-dimensional spectra can produce a
\textbf{non-physical} spectrum (a peak pattern no instrument would ever record), so pair SMOTE with \emph{curation}
and domain augmentation ($m/z$ jitter), or simply \emph{oversample real} spectra.}

\vspace{4pt}\hrule\vspace{3pt}
{\small\itshape \ifsolution Answer key --- other defensible wording earns credit if the numbers match.%
\else Keep this sheet --- it's your checklist for any rare-class assay in your own lab.\fi}

\end{document}
"
%  (or, to flip by hand, change \solutionfalse to \solutiontrue below.)
% =====================================================================
\documentclass[11pt]{article}

\newif\ifsolution
\ifdefined\showsolutions \solutiontrue \else \solutionfalse \fi
% \solutiontrue   % <-- uncomment to hard-code the ANSWER KEY build

\usepackage[letterpaper, margin=0.6in, top=0.7in, bottom=0.5in]{geometry}
\usepackage{amsmath, amssymb}
\usepackage[table]{xcolor}
\usepackage{array}
\usepackage{fancyhdr}

\pagestyle{fancy}
\fancyhf{}
\fancyhead[L]{\small\textbf{MSACL DS301 --- Deep Learning}}
\fancyhead[R]{\small\textbf{Quiz 10: Imbalanced Data\ifsolution\ \textmd{(Answer Key)}\fi}}
\fancyfoot[C]{\thepage}
\renewcommand{\headrulewidth}{0.4pt}

\newcommand{\blk}[1]{\underline{\hspace{#1}}}
% Reveal-or-blank: shows the answer in the key, a blank rule on the student sheet.
\newcommand{\rb}[2]{\ifsolution\textbf{#1}\else\blk{#2}\fi}
% Boxed answer (key only) and blank writing space (student only).
\newcommand{\keybox}[1]{\ifsolution\par\vspace{2pt}%
  \fbox{\parbox{0.965\textwidth}{\small\textbf{Answer.}\; #1}}\par\fi}
\newcommand{\work}[1]{\ifsolution\else\vspace{#1}\fi}

% ---- course palette (numeric grids) --------------------------------------
\definecolor{ink}{HTML}{1B2025}
\definecolor{teal}{HTML}{0E7C7B}
\definecolor{tealsoft}{HTML}{E3F0EF}
\definecolor{amber}{HTML}{E09E2F}
\definecolor{ambersoft}{HTML}{FBF0DC}
\definecolor{redd}{HTML}{C4534F}
\definecolor{redsoft}{HTML}{F7E4E3}
\definecolor{inksoft}{HTML}{3D444C}
\definecolor{muted}{HTML}{8A9099}
\definecolor{hair}{HTML}{D9D9D2}
% square, centred cells for the confusion matrix
\newcolumntype{G}{>{\centering\arraybackslash}p{28mm}}
\newcommand{\gr}{\renewcommand{\arraystretch}{1.7}}

\setlength{\parindent}{0pt}
\setlength{\parskip}{3pt}
\setlength{\abovedisplayskip}{5pt}
\setlength{\belowdisplayskip}{5pt}

\begin{document}

\begin{center}
{\Large \textbf{Quiz 10 --- Imbalanced Data}}\\[2pt]
{\normalsize Compute precision / recall / specificity, choose AUROC vs.\ AUPRC, match a treatment, and interpolate a SMOTE point}
\end{center}

\vspace{-2pt}
\begin{center}
\fbox{\parbox{0.92\textwidth}{\small
\textbf{Instructions:}\; Answer all four questions --- every part is answerable from the formulas and pictures on
the slides. A rapid resistance screen was run on \textbf{500 isolates} (positive $=$ \emph{resistant}); Q1--Q2 use
its confusion matrix. Compute each rate exactly; no calculators needed.%
\ifsolution\else\\[4pt]\textbf{Name:}\ \blk{6cm}\hfill\textbf{Date:}\ \blk{3.5cm}\fi
}}
\end{center}

\vspace{-1pt}
\begin{center}\small
\textbf{Formulas (from the slides):}\;
$\text{recall (sensitivity)} = \dfrac{TP}{TP+FN}$ \quad
$\text{specificity} = \dfrac{TN}{TN+FP}$ \quad
$\text{precision} = \dfrac{TP}{TP+FP}$ \quad
$\text{FPR} = \dfrac{FP}{FP+TN}$
\end{center}

\vspace{-2pt}\hrule\vspace{5pt}

% =========================== Q1 ===========================
\textbf{1. Three rates by hand.}\; {\small The screen's confusion matrix (rows $=$ the model's call, columns $=$
truth):}

\vspace{4pt}
\begin{center}
\gr\begin{tabular}{r|G|G|}
\multicolumn{1}{r}{} & \multicolumn{1}{c}{\small truth: $+$ (resistant)} & \multicolumn{1}{c}{\small truth: $-$ (susceptible)}\\\cline{2-3}
\small call: $+$ & \cellcolor{tealsoft}TP $=$ 16 & \cellcolor{ambersoft}FP $=$ 32 \\\cline{2-3}
\small call: $-$ & \cellcolor{redsoft}FN $=$ 4 & \cellcolor{tealsoft}TN $=$ 448 \\\cline{2-3}
\end{tabular}\\[3pt]
{\footnotesize 20 resistant \;$\cdot$\; 480 susceptible \;$\cdot$\; 500 total}
\end{center}

\vspace{2pt}
\textbf{(a)} recall (sensitivity) $=$ \rb{$16/20 = 0.80$ \,(80\%)}{3.4cm}
\work{0.3cm}

\textbf{(b)} specificity $=$ \rb{$448/480 = 0.93$ \,(93\%)}{3.4cm}
\work{0.3cm}

\textbf{(c)} precision $=$ \rb{$16/48 = 0.33$ \,(33\%)}{3.4cm}
\work{0.3cm}

\textbf{(d)} A model that ALWAYS calls ``susceptible'' scores accuracy $= 480/500 = 96\%$. The lab must not miss
resistant isolates. \emph{Which rate(s) do you trust here, and why is accuracy misleading?}

\work{1.0cm}
\ifsolution\else\blk{\linewidth}\par\vspace{2pt}\blk{\linewidth}\fi

\keybox{(a) recall $= 16/(16+4) = 16/20 = \mathbf{0.80}$ (\textbf{80\%}). (b) specificity $= 448/(448+32) = 448/480 =
\mathbf{0.93}$ (\textbf{93\%}). (c) precision $= 16/(16+32) = 16/48 = \mathbf{0.33}$ (\textbf{33\%}). Accuracy $=
(16+448)/500 = 464/500 = \mathbf{0.93}$ (93\%). (d) Trust \textbf{recall} (catch the resistant) reported \emph{with}
\textbf{precision} (only 1 flag in 3 is real). Positives are rare (20 of 500), so accuracy is dominated by the 480
susceptibles: an always-``susceptible'' model scores \textbf{96\%} yet catches \textbf{0 of 20} resistant. Accuracy
hides both the missed cases and the false alarms.}

\vspace{2pt}\hrule\vspace{5pt}

% =========================== Q2 ===========================
\textbf{2. AUROC or AUPRC?}\; {\small Same screen. Positives are \textbf{rare: 20 of 500 $= 4\%$ prevalence.} Plot
this one operating point on each curve, then choose what to report across all thresholds.}

\vspace{2pt}
\textbf{(a)} ROC point: FPR $=\ \dfrac{32}{32+448} =\ $ \rb{$32/480 = 0.067$}{2.8cm}, \; TPR (recall) $=\ $
\rb{$0.80$}{2.0cm}\; --- does it hug the top-left corner? \rb{yes (looks great)}{3.0cm}
\work{0.3cm}

\textbf{(b)} PR point: recall $=\ $ \rb{$0.80$}{2.0cm}, \; precision $=\ $ \rb{$0.33$}{2.0cm};\;
the PR chance baseline sits at the prevalence $=\ $ \rb{$20/500 = 0.04$}{2.8cm}
\work{0.3cm}

\textbf{(c)} To report this screen across \emph{all thresholds}, would you trust the \textbf{AUROC} or the
\textbf{AUPRC / PR curve}, and why?

\work{1.0cm}
\ifsolution\else\blk{\linewidth}\par\vspace{2pt}\blk{\linewidth}\fi

\keybox{(a) FPR $= 32/(32+448) = 32/480 = \mathbf{0.067}$; TPR $= \mathbf{0.80}$ --- the point $(0.067,\,0.80)$
\textbf{hugs the top-left}, so the ROC/AUROC \emph{looks excellent}. (b) PR point $=(\text{recall }\mathbf{0.80},\,
\text{precision }\mathbf{0.33})$; PR baseline $=$ prevalence $= 20/500 = \mathbf{0.04}$ (ROC's baseline is $0.5$).
(c) Trust the \textbf{AUPRC / PR curve} (reported with recall). At 4\% prevalence the 480 true-negatives dominate, so a
tiny FPR keeps AUROC flatteringly high even though precision is only \textbf{0.33} (2 of 3 flags false). AUROC is
invariant to the class ratio; AUPRC collapses toward the prevalence, so it stays honest as recall grows
(Saito \& Rehmsmeier 2015).}

\vspace{2pt}\hrule\vspace{5pt}

% =========================== Q3 ===========================
\textbf{3. Match the treatment.}\; {\small Match each situation to \emph{one} move from the menu:
\textbf{curation} \;$\cdot$\; \textbf{class weights / SMOTE / oversample} \;$\cdot$\; \textbf{focal loss}
\;$\cdot$\; \textbf{curriculum}. Each is used once.}

\vspace{2pt}
\textbf{(a)} Junk runs and mislabeled spectra contaminate the training set --- and the same isolate
appears three times. $\rightarrow$\ \rb{curation}{5cm}
\work{0.25cm}
\textbf{(b)} Only 41 real resistant spectra; the model just predicts ``susceptible''. $\rightarrow$\
\rb{class weights / SMOTE / oversample}{5cm}
\work{0.25cm}
\textbf{(c)} Easy cases are fine, but borderline isolates are missed. $\rightarrow$\ \rb{focal loss}{5cm}
\work{0.25cm}
\textbf{(d)} Training is unstable when very easy and very hard batches alternate. $\rightarrow$\
\rb{curriculum}{5cm}
\work{0.35cm}

\textbf{(e)} \emph{Which of these four adds NEW resistant signal to the dataset?}

\work{0.9cm}
\ifsolution\else\blk{\linewidth}\par\fi

\keybox{(a) \textbf{curation} --- clean the labels and drop the junk before anything fancier, then
\emph{deduplicate}: embed every spectrum (the CNN's last hidden layer) and treat $\cos \geq 0.99$ as the
same isolate, keeping one --- repeats split across train and test are leakage. (b) \textbf{class weights / SMOTE /
oversample} --- make the rare class count. (c) \textbf{focal loss} --- up-weight the hard, low-confidence cases.
(d) \textbf{curriculum} --- order easy$\to$hard, or class-balanced$\to$true prevalence. (e) \textbf{None of them.}
Every move re-uses the same 41 resistant spectra (reweighting shouts them louder, SMOTE interpolates between them,
focal loss focuses on them, curriculum reorders them) --- \emph{no new biological information enters}. Only
\textbf{collecting more real resistant isolates} adds signal.}

\vspace{2pt}\hrule\vspace{5pt}

% =========================== Q4 ===========================
\textbf{4. SMOTE by hand.}\; {\small SMOTE makes a new minority point on the segment between two real ones:
\; $\text{synthetic} = A + \lambda\,(B-A)$. Use two $m/z$-bin intensities $A=(2,\,6)$ and $B=(4,\,10)$.}

\vspace{2pt}
\textbf{(a)} With $\lambda = 0.25$:\; synthetic $=\ (\,2 + 0.25\,(4-2),\ \ 6 + 0.25\,(10-6)\,) =\ $
\rb{$(2.5,\ 7)$}{3cm}
\work{0.4cm}

\textbf{(b)} \emph{One sentence: why is raw-spectrum SMOTE risky in 6{,}000 dimensions?}

\work{1.0cm}
\ifsolution\else\blk{\linewidth}\par\vspace{2pt}\blk{\linewidth}\fi

\keybox{(a) synthetic $= (2 + 0.25\cdot 2,\ 6 + 0.25\cdot 4) = (2 + 0.5,\ 6 + 1) = \mathbf{(2.5,\ 7)}$ --- a quarter
of the way from $A$ toward $B$. (b) A straight-line blend of two real 6{,}000-dimensional spectra can produce a
\textbf{non-physical} spectrum (a peak pattern no instrument would ever record), so pair SMOTE with \emph{curation}
and domain augmentation ($m/z$ jitter), or simply \emph{oversample real} spectra.}

\vspace{4pt}\hrule\vspace{3pt}
{\small\itshape \ifsolution Answer key --- other defensible wording earns credit if the numbers match.%
\else Keep this sheet --- it's your checklist for any rare-class assay in your own lab.\fi}

\end{document}
"
%  (or, to flip by hand, change \solutionfalse to \solutiontrue below.)
% =====================================================================
\documentclass[11pt]{article}

\newif\ifsolution
\ifdefined\showsolutions \solutiontrue \else \solutionfalse \fi
% \solutiontrue   % <-- uncomment to hard-code the ANSWER KEY build

\usepackage[letterpaper, margin=0.6in, top=0.7in, bottom=0.5in]{geometry}
\usepackage{amsmath, amssymb}
\usepackage[table]{xcolor}
\usepackage{array}
\usepackage{fancyhdr}

\pagestyle{fancy}
\fancyhf{}
\fancyhead[L]{\small\textbf{MSACL DS301 --- Deep Learning}}
\fancyhead[R]{\small\textbf{Quiz 10: Imbalanced Data\ifsolution\ \textmd{(Answer Key)}\fi}}
\fancyfoot[C]{\thepage}
\renewcommand{\headrulewidth}{0.4pt}

\newcommand{\blk}[1]{\underline{\hspace{#1}}}
% Reveal-or-blank: shows the answer in the key, a blank rule on the student sheet.
\newcommand{\rb}[2]{\ifsolution\textbf{#1}\else\blk{#2}\fi}
% Boxed answer (key only) and blank writing space (student only).
\newcommand{\keybox}[1]{\ifsolution\par\vspace{2pt}%
  \fbox{\parbox{0.965\textwidth}{\small\textbf{Answer.}\; #1}}\par\fi}
\newcommand{\work}[1]{\ifsolution\else\vspace{#1}\fi}

% ---- course palette (numeric grids) --------------------------------------
\definecolor{ink}{HTML}{1B2025}
\definecolor{teal}{HTML}{0E7C7B}
\definecolor{tealsoft}{HTML}{E3F0EF}
\definecolor{amber}{HTML}{E09E2F}
\definecolor{ambersoft}{HTML}{FBF0DC}
\definecolor{redd}{HTML}{C4534F}
\definecolor{redsoft}{HTML}{F7E4E3}
\definecolor{inksoft}{HTML}{3D444C}
\definecolor{muted}{HTML}{8A9099}
\definecolor{hair}{HTML}{D9D9D2}
% square, centred cells for the confusion matrix
\newcolumntype{G}{>{\centering\arraybackslash}p{28mm}}
\newcommand{\gr}{\renewcommand{\arraystretch}{1.7}}

\setlength{\parindent}{0pt}
\setlength{\parskip}{3pt}
\setlength{\abovedisplayskip}{5pt}
\setlength{\belowdisplayskip}{5pt}

\begin{document}

\begin{center}
{\Large \textbf{Quiz 10 --- Imbalanced Data}}\\[2pt]
{\normalsize Compute precision / recall / specificity, choose AUROC vs.\ AUPRC, match a treatment, and interpolate a SMOTE point}
\end{center}

\vspace{-2pt}
\begin{center}
\fbox{\parbox{0.92\textwidth}{\small
\textbf{Instructions:}\; Answer all four questions --- every part is answerable from the formulas and pictures on
the slides. A rapid resistance screen was run on \textbf{500 isolates} (positive $=$ \emph{resistant}); Q1--Q2 use
its confusion matrix. Compute each rate exactly; no calculators needed.%
\ifsolution\else\\[4pt]\textbf{Name:}\ \blk{6cm}\hfill\textbf{Date:}\ \blk{3.5cm}\fi
}}
\end{center}

\vspace{-1pt}
\begin{center}\small
\textbf{Formulas (from the slides):}\;
$\text{recall (sensitivity)} = \dfrac{TP}{TP+FN}$ \quad
$\text{specificity} = \dfrac{TN}{TN+FP}$ \quad
$\text{precision} = \dfrac{TP}{TP+FP}$ \quad
$\text{FPR} = \dfrac{FP}{FP+TN}$
\end{center}

\vspace{-2pt}\hrule\vspace{5pt}

% =========================== Q1 ===========================
\textbf{1. Three rates by hand.}\; {\small The screen's confusion matrix (rows $=$ the model's call, columns $=$
truth):}

\vspace{4pt}
\begin{center}
\gr\begin{tabular}{r|G|G|}
\multicolumn{1}{r}{} & \multicolumn{1}{c}{\small truth: $+$ (resistant)} & \multicolumn{1}{c}{\small truth: $-$ (susceptible)}\\\cline{2-3}
\small call: $+$ & \cellcolor{tealsoft}TP $=$ 16 & \cellcolor{ambersoft}FP $=$ 32 \\\cline{2-3}
\small call: $-$ & \cellcolor{redsoft}FN $=$ 4 & \cellcolor{tealsoft}TN $=$ 448 \\\cline{2-3}
\end{tabular}\\[3pt]
{\footnotesize 20 resistant \;$\cdot$\; 480 susceptible \;$\cdot$\; 500 total}
\end{center}

\vspace{2pt}
\textbf{(a)} recall (sensitivity) $=$ \rb{$16/20 = 0.80$ \,(80\%)}{3.4cm}
\work{0.3cm}

\textbf{(b)} specificity $=$ \rb{$448/480 = 0.93$ \,(93\%)}{3.4cm}
\work{0.3cm}

\textbf{(c)} precision $=$ \rb{$16/48 = 0.33$ \,(33\%)}{3.4cm}
\work{0.3cm}

\textbf{(d)} A model that ALWAYS calls ``susceptible'' scores accuracy $= 480/500 = 96\%$. The lab must not miss
resistant isolates. \emph{Which rate(s) do you trust here, and why is accuracy misleading?}

\work{1.0cm}
\ifsolution\else\blk{\linewidth}\par\vspace{2pt}\blk{\linewidth}\fi

\keybox{(a) recall $= 16/(16+4) = 16/20 = \mathbf{0.80}$ (\textbf{80\%}). (b) specificity $= 448/(448+32) = 448/480 =
\mathbf{0.93}$ (\textbf{93\%}). (c) precision $= 16/(16+32) = 16/48 = \mathbf{0.33}$ (\textbf{33\%}). Accuracy $=
(16+448)/500 = 464/500 = \mathbf{0.93}$ (93\%). (d) Trust \textbf{recall} (catch the resistant) reported \emph{with}
\textbf{precision} (only 1 flag in 3 is real). Positives are rare (20 of 500), so accuracy is dominated by the 480
susceptibles: an always-``susceptible'' model scores \textbf{96\%} yet catches \textbf{0 of 20} resistant. Accuracy
hides both the missed cases and the false alarms.}

\vspace{2pt}\hrule\vspace{5pt}

% =========================== Q2 ===========================
\textbf{2. AUROC or AUPRC?}\; {\small Same screen. Positives are \textbf{rare: 20 of 500 $= 4\%$ prevalence.} Plot
this one operating point on each curve, then choose what to report across all thresholds.}

\vspace{2pt}
\textbf{(a)} ROC point: FPR $=\ \dfrac{32}{32+448} =\ $ \rb{$32/480 = 0.067$}{2.8cm}, \; TPR (recall) $=\ $
\rb{$0.80$}{2.0cm}\; --- does it hug the top-left corner? \rb{yes (looks great)}{3.0cm}
\work{0.3cm}

\textbf{(b)} PR point: recall $=\ $ \rb{$0.80$}{2.0cm}, \; precision $=\ $ \rb{$0.33$}{2.0cm};\;
the PR chance baseline sits at the prevalence $=\ $ \rb{$20/500 = 0.04$}{2.8cm}
\work{0.3cm}

\textbf{(c)} To report this screen across \emph{all thresholds}, would you trust the \textbf{AUROC} or the
\textbf{AUPRC / PR curve}, and why?

\work{1.0cm}
\ifsolution\else\blk{\linewidth}\par\vspace{2pt}\blk{\linewidth}\fi

\keybox{(a) FPR $= 32/(32+448) = 32/480 = \mathbf{0.067}$; TPR $= \mathbf{0.80}$ --- the point $(0.067,\,0.80)$
\textbf{hugs the top-left}, so the ROC/AUROC \emph{looks excellent}. (b) PR point $=(\text{recall }\mathbf{0.80},\,
\text{precision }\mathbf{0.33})$; PR baseline $=$ prevalence $= 20/500 = \mathbf{0.04}$ (ROC's baseline is $0.5$).
(c) Trust the \textbf{AUPRC / PR curve} (reported with recall). At 4\% prevalence the 480 true-negatives dominate, so a
tiny FPR keeps AUROC flatteringly high even though precision is only \textbf{0.33} (2 of 3 flags false). AUROC is
invariant to the class ratio; AUPRC collapses toward the prevalence, so it stays honest as recall grows
(Saito \& Rehmsmeier 2015).}

\vspace{2pt}\hrule\vspace{5pt}

% =========================== Q3 ===========================
\textbf{3. Match the treatment.}\; {\small Match each situation to \emph{one} move from the menu:
\textbf{curation} \;$\cdot$\; \textbf{class weights / SMOTE / oversample} \;$\cdot$\; \textbf{focal loss}
\;$\cdot$\; \textbf{curriculum}. Each is used once.}

\vspace{2pt}
\textbf{(a)} Junk runs and mislabeled spectra contaminate the training set --- and the same isolate
appears three times. $\rightarrow$\ \rb{curation}{5cm}
\work{0.25cm}
\textbf{(b)} Only 41 real resistant spectra; the model just predicts ``susceptible''. $\rightarrow$\
\rb{class weights / SMOTE / oversample}{5cm}
\work{0.25cm}
\textbf{(c)} Easy cases are fine, but borderline isolates are missed. $\rightarrow$\ \rb{focal loss}{5cm}
\work{0.25cm}
\textbf{(d)} Training is unstable when very easy and very hard batches alternate. $\rightarrow$\
\rb{curriculum}{5cm}
\work{0.35cm}

\textbf{(e)} \emph{Which of these four adds NEW resistant signal to the dataset?}

\work{0.9cm}
\ifsolution\else\blk{\linewidth}\par\fi

\keybox{(a) \textbf{curation} --- clean the labels and drop the junk before anything fancier, then
\emph{deduplicate}: embed every spectrum (the CNN's last hidden layer) and treat $\cos \geq 0.99$ as the
same isolate, keeping one --- repeats split across train and test are leakage. (b) \textbf{class weights / SMOTE /
oversample} --- make the rare class count. (c) \textbf{focal loss} --- up-weight the hard, low-confidence cases.
(d) \textbf{curriculum} --- order easy$\to$hard, or class-balanced$\to$true prevalence. (e) \textbf{None of them.}
Every move re-uses the same 41 resistant spectra (reweighting shouts them louder, SMOTE interpolates between them,
focal loss focuses on them, curriculum reorders them) --- \emph{no new biological information enters}. Only
\textbf{collecting more real resistant isolates} adds signal.}

\vspace{2pt}\hrule\vspace{5pt}

% =========================== Q4 ===========================
\textbf{4. SMOTE by hand.}\; {\small SMOTE makes a new minority point on the segment between two real ones:
\; $\text{synthetic} = A + \lambda\,(B-A)$. Use two $m/z$-bin intensities $A=(2,\,6)$ and $B=(4,\,10)$.}

\vspace{2pt}
\textbf{(a)} With $\lambda = 0.25$:\; synthetic $=\ (\,2 + 0.25\,(4-2),\ \ 6 + 0.25\,(10-6)\,) =\ $
\rb{$(2.5,\ 7)$}{3cm}
\work{0.4cm}

\textbf{(b)} \emph{One sentence: why is raw-spectrum SMOTE risky in 6{,}000 dimensions?}

\work{1.0cm}
\ifsolution\else\blk{\linewidth}\par\vspace{2pt}\blk{\linewidth}\fi

\keybox{(a) synthetic $= (2 + 0.25\cdot 2,\ 6 + 0.25\cdot 4) = (2 + 0.5,\ 6 + 1) = \mathbf{(2.5,\ 7)}$ --- a quarter
of the way from $A$ toward $B$. (b) A straight-line blend of two real 6{,}000-dimensional spectra can produce a
\textbf{non-physical} spectrum (a peak pattern no instrument would ever record), so pair SMOTE with \emph{curation}
and domain augmentation ($m/z$ jitter), or simply \emph{oversample real} spectra.}

\vspace{4pt}\hrule\vspace{3pt}
{\small\itshape \ifsolution Answer key --- other defensible wording earns credit if the numbers match.%
\else Keep this sheet --- it's your checklist for any rare-class assay in your own lab.\fi}

\end{document}
